\section{Мета роботи та завдання}

\subsection{Мета лабораторної роботи}

Метою лабораторної роботи є формування стійких практичних навичок роботи 
з обліковими структурами даних та стандартними колекціями в мові програмування 
Java. У процесі виконання роботи необхідно:
\begin{itemize}
    \item 
    Ознайомитися з основними архітектурними видами колекцій Java (\texttt{List}, 
    \texttt{Set}, \texttt{Map});
    \item 
    Навчитися проектувати динамічні списки, множини та асоціативні 
    таблиці відображень;
    \item 
    Опанувати базові операції додавання, пошуку, видалення та ітераційного 
    перебору елементів колекцій;
    \item 
    Практично застосувати механізм узагальнень (\textit{Generics}) для 
    забезпечення суворої типобезпеки на етапі компіляції додатка.
\end{itemize}

\subsection{Загальне технічне завдання (Практична частина)}

Програма досліджень вимагає розробки консольного додатка, що демонструє 
базові операції над трьома типами колекцій:
\begin{enumerate}
    \item \textbf{Робота зі списками (\texttt{List<Student>}):} 
    Створення динамічного списку об'єктів користувацького класу, виведення 
    елементів за допомогою циклу \texttt{for-each}, розрахунок середнього 
    арифметичного балу групи за допомогою методу \texttt{size()} та пошук 
    лідера за успішністю.
    \item \textbf{Робота з множинами (\texttt{Set<String>}):} 
    Організація сховища унікальних навчальних дисциплін на 
    базі \texttt{HashSet} з демонстрацією автоматичного відсікання дублікатів 
    при спробі повторного додавання елементів.
    \item \textbf{Робота з відображеннями (\texttt{Map<Integer, String>}):} 
    Реалізація індексної таблиці відповідності між унікальним кодом (ідентифікатором) 
    студента та його ПІБ за допомогою \texttt{HashMap}. Демонстрація пошуку за ключем, 
    перевірки наявності індексу та виведення записів через \texttt{entrySet()}.
\end{enumerate}

\subsection{Індивідуальне завдання (Варіант №7)}

Розробити консольну програму для обліку фільмів. Створити користувацький 
клас \texttt{Movie}, який містить назву, жанр, рік випуску та поточний рейтинг 
кінострічки. 

Для організації збереження даних необхідно використати динамічну 
колекцію \texttt{List<Movie>}. Програма повинна вивести повний список фільмів, 
розробити лінійний алгоритм визначення фільму з найвищим рейтингом у списку, а 
також автоматично сформувати та вивести множину унікальних жанрів кінематографа 
за допомогою структури \texttt{Set<String>}.

\subsection{Оформити звіт і відповісти на контрольні питання.}


